% ======================================================================
% 10. SYNTHESIS  (budget: 1.0 page)
% Job of the chapter: answer the question the title implies -- so when does
% this still work? A decision chain, not a summary of the paper.
% ======================================================================
\section{When These Attacks Actually Work}
\label{sec:synthesis}

\subsection{The precondition chain as a decision procedure}
% CLAIM (0.4): five questions in order, each one able to end the attack:
%   1. Is NTLM accepted anywhere on the path?        -> PtH viable
%   2. Is RC4 enabled for the account?               -> PtK viable
%   3. Is the hash reachable on a host we already own? -> both viable
%   4. Is extraction blocked (Credential Guard/PPL)? -> extraction dead, live
%                                                       session reuse is not
%   5. Is the account authorised on the target?      -> the real bottleneck
% Most real intrusions die at 5, and almost nobody writes that down.

\subsection{Where it fails in a modern estate}
% CLAIM (0.3): concrete failure modes we observed plus the ones we can argue
% from the protocol: token filtering, NTLM disabled per service, AES-only,
% Protected Users, tiered admin with no overlapping local admins.

\subsection{What survives the deprecation}
% CLAIM (0.3): NTLM deprecation and RC4 removal shrink this space over the
% next few years, but the class of attack survives as long as any protocol
% authenticates on derived key material -- which is exactly why ticket and
% certificate abuse are the same argument in different clothes.
